\pg{425} pela sua \emph{causa efficiente}. Ora a existencia de \textsc{Deus} não se prova pela sua \emph{causa efficiente}; porque \textsc{Deus} é a Causa primeira, e por isso absolutamente independente de toda e qualquer causa. Logo a existencia de \textsc{Deus} não se demonstra \emph{a-priori}.

\textbf{9. A existencia de Deus não se demonstra a-simultaneo.} --- A demonstração é \emph{a-simultaneo}, quando a existencia de um ente se prova pela analyse da idea ou noção do proprio ente. Ora este processo, applicado á existencia de \textsc{Deus}, é vicioso; visto que, para uma coisa existir na realidade, não basta que exista no pensamento, pois deduzir da ordem \emph{ideal} a ordem \emph{real} é tirar uma consequencia mais extensa que as premissas. Logo a existencia de \textsc{Deus} não se demonstra \emph{a-simultaneo}\footnote{Esta proposição --- \textsc{DEus} \emph{existe} --- é \emph{immediatamente} evidente \emph{em si} mesma, mas não é \emph{immediatamente} evidente \emph{em relação a} nós.\par a) E' \emph{immediatamente} evidente \emph{em si mesma}. Uma proposição é \emph{immediatamente} evidente \emph{em si mesma}, quando o predicado convém \emph{essencialmente} ao sujeito. Ora a \emph{existencia} convém \emph{essencialmente} a \textsc{DEus}; porque \textsc{DEus} é o ser \emph{necessario}, e não póde deixar de existir. Logo a proposição --- \textsc{DEus} \emph{existe} --- é \emph{immediatamente} evidente \emph{em si mesma}.\par \emph{b}) Não é \emph{immediatamente} evidente \emph{em relação a nós}. Uma proposição é \emph{immediatamente} evidente \emph{em relação a nós}, quando directamente e sem necessidade de raciocinio conhecemos a essencia do sujeito e do predicado; assim a proposição --- \emph{o todo é maior que a parte} --- é \emph{immediatamente} evidente \emph{em relação a nós}, pois basta enunciar as palavras \emph{todo} e \emph{parte}, para logo percebermos que o \emph{todo} é maior que a \emph{parte}. Ora a essencia de \textsc{DEus}, não sendo por nós conhecida senão depois de um raciocinio, e ainda assim muito imperfeitamente, não se nos manifesta logo á primeira vista como exigindo a existencia. Logo a proposição --- \textsc{DEus} \emph{existe} --- não é \emph{immediatamente} evidente \emph{em relação a nós}. --- Cf. S. Thom., --- \emph{Sum}. \emph{Th}., p. I, q. 2, a. 1; --- c. \emph{Gentes}., l. I, cc. 10, 11; 1 \emph{Sent}. Dist. 3, q. 1, a. 2; \emph{De verit}., q. 10, a. 12.\par Os mais celebres argumentos \emph{a-simultaneo} em favor da existencia de \textsc{DEus} são os de \emph{Santo Anselmo}, \emph{Descartes} e \emph{Leibnitz}.\par a) \emph{Santo Anselmo} raciocina d'este modo: Uma proposição é \emph{immediatamente} evidente \emph{em relação a nós}, quando a sua verdade se nos manifesta pela simples enunciação dos termos. Ora a verdade d'esta proposição --- \textsc{DEus} \emph{existe} --- torna-se-nos manifesta pela simples enunciação dos termos. Com effeito, quando se profere o nome de \textsc{DEus}, immediatamente entendemos que \textsc{DEus} é o Ser maior que se póde imaginar, e que por isso deve existir na \emph{realidade}; porque, se existisse unicamente na intelligencia, não seria o Ser maior que se póde imaginar, pois superior a Elle seria o ser, que existisse na intelligencia e tambem na \emph{realidade}.\par Critica. Ainda que todos admittissem que \textsc{DEus} é o ser maior que se póde imaginar, todavia d'ahi só poderia concluir-se que \textsc{DEus} existe na \emph{intelligencia}, mas não na \emph{realidade}; porque, como dissemos, deduzir da ordem \emph{ideal} a ordem \emph{real} é tirar uma consequencia, que não está incluida nas premissas\pgn{426}. --- Para provar que \textsc{DEus} existe tambem na \emph{realidade}, deve demonstrar-se que existe na \emph{realidade} um ser, que é o maior que se póde imaginar. --- Cf. S. Thom., c. \emph{Gentes}, l. I, c. 11.\par b) \emph{Descartes} diz: Todo o attribuito, que está contido na idea \emph{clara} e \emph{distincta} de um sujeito, deve affirmar-se do mesmo sujeito. Ora na idea \emph{clara} e \emph{distincta} de \textsc{DEus} está contida a existencia \emph{real}; porque \textsc{DEus} é infinitamente perfeito, e a existencia \emph{real} é uma perfeição. Logo deve affirmar-se de \textsc{DEus} a existencia \emph{real}, isto é, \textsc{DEus} existe \emph{realmente}.\par Critica. O argumento de \emph{Descartes} tem o mesmo defeito que o de \emph{Santo Anselmo}; pois passa da ordem \emph{ideal} para a \emph{real}. Na verdade, tudo o que está contido na idea \emph{clara} e \emph{distincta} de um sujeito deve attribuir-se ao sujeito na mesma \emph{ordem}, em que está o sujeito; de modo que, se o sujeito estiver na ordem \emph{real}, o predicado tambem estará na ordem \emph{real}, --- mas, se o sujeito estiver na ordem \emph{ideal}, o predicado tambem existirá necessariamente na ordem \emph{ideal}. Logo na \emph{idea} de \textsc{DEus} está contida a existencia \emph{ideal}, e não a \emph{real}. --- Lembramos que, para \emph{Descartes}, a idea do \emph{infinito} é \emph{innata} (o que é falso), e que a verdade da proposição maior do seu syllogismo se baséa na veracidade do proprio \textsc{DEus} (o que constitue um \emph{circulo vicioso}).\par c) \emph{Leibnitz} apresenta este argumento: O ente infinitamente perfeito é possivel. Se é possivel, existe na \emph{realidade}; porque a existencia \emph{real}, por ser uma perfeição, pertence necessariamente á essencia do Ente infinitamente perfeito. Logo o Ente infinitamente perfeito existe na \emph{realidade}.\par Critica. Este argumento é inconcludente pelo mesmo defeito, que apontámos nos dois antecedentes. Do facto de \emph{concebermos}, como \emph{possivel}, o Ente infinitamente perfeito, cuja essencia exige a existencia, só resulta que não podemos \emph{conceber}, como \emph{possivel}, esse ente, sem que o \emph{concebamos} como existente; mas não se segue que a realidade deva corresponder ao que \emph{concebemos} na intelligencia.}.

\textbf{10. A existencia de Deus demonstra-se a-posteriori.} --- A demonstração \emph{a-posteriori} faz-se, quando, pela existencia do effeito,